\section{Introducción al curso: de qué trata esta clase}

Esta clase es la introducción de apertura de «Large Language Model Agents» de Berkeley, en la que Dawn Song presenta el enfoque del curso, el equipo docente y los objetivos de aprendizaje. Ella sitúa a los Agentes en el contexto amplio de la rápida evolución de los LLM: los modelos de lenguaje ya son suficientemente potentes, pero para resolver tareas reales es necesario que el modelo forme un ciclo cerrado con el entorno, las herramientas, la memoria y la retroalimentación.

\begin{figure}[H]
\centering
\includegraphics[page=2,width=0.9\textwidth]{slides.pdf}
\caption{Equipo docente del curso y roles de colaboración. La introducción primero deja claros el cuerpo docente, los ayudantes y los recursos de apoyo del curso.}
\end{figure}

\begin{knowledgebox}{El punto de partida de este curso}
El curso no trata a los Agentes como un conjunto de técnicas de un solo prompt, sino que los define como un sistema completo: el LLM se encarga del razonamiento y la planificación, mientras que los módulos externos se encargan del uso de herramientas, la memoria, la recuperación y la interacción con el entorno. Solo al encadenar estas capacidades puede el LLM pasar de «responder preguntas» a «completar tareas».
\end{knowledgebox}

\subsection{Por qué los modelos grandes avanzan de manera natural hacia los Agentes}

La introducción señala que, en los últimos dos años, la capacidad de los LLM ha mejorado con gran rapidez, pero el modelo de entrada de puro texto a salida de puro texto pronto alcanzó su límite. Las tareas del mundo real por lo general no son preguntas y respuestas de una sola vez, sino que requieren:
\begin{itemize}
    \item interactuar con el entorno, leer páginas, archivos, bases de datos o el estado de dispositivos;
    \item conservar el contexto de la tarea y la memoria a largo plazo;
    \item reintentar tras un fallo, ajustar el plan y corregir el comportamiento a partir de la retroalimentación;
    \item invocar herramientas externas para suplir las carencias del modelo en cómputo, recuperación y ejecución.
\end{itemize}

\begin{importantbox}{El salto clave de LLM a Agente}
Una vez que el modelo cuenta con memoria, uso de herramientas, recuperación y un ciclo de retroalimentación, su unidad de comportamiento deja de ser «responder una frase» para convertirse en «hacer avanzar una tarea a lo largo de una interacción de varios pasos». Esta es también la línea divisoria fundamental entre la investigación de Agentes y las aplicaciones comunes de chatbot.
\end{importantbox}

\subsection{Resumen de la sección}

La introducción del curso deja claro el planteamiento desde el inicio: un Agente no consiste en ponerle otra capa de interfaz de usuario a un LLM, sino en colocar el modelo dentro de un marco de ejecución que observa, planifica, actúa y corrige de manera continua.
